\documentclass[10pt,a4paper]{amsart}
\usepackage[margin=19mm]{geometry}
\usepackage{amsmath,amssymb,mathtools}
\usepackage[scr=rsfs,cal=euler]{mathalfa}
\usepackage{xfrac}
\usepackage[unicode,hidelinks,bookmarks=false]{hyperref}
\newcommand{\ie}{i.e.\ }
\newcommand{\cf}{cf.\ }
\newcommand{\resp}{respectively\ }
\newcommand{\Subsection}{\subsection}
\providecommand{\nobreakdash}{\nobreak\hbox{-}}
\setlength{\emergencystretch}{2em}
\multlinegap0pt
\usepackage{amssymb}
\usepackage{bm}
\newtheorem{theorem}{Theorem}
\def\AA{A}
\def\BDst{\dd_{\!\!}\Bpdst}
\def\Bpdst{{}_{\,}{\reflectbox{\mbox{$\Bpstd$}}}}
\def\Bpstd{{}^{\phantom{J}}_{\raisebox{\depth}{\scalebox{.75}[-.75]{\mbox{${\bm
            \Delta}$}}}}}
\def\MMM{\mathscr{M}}
\def\Vect{\mathfrak{X}}
\def\al{\alpha}
\def\be{\beta}
\def\catfp{{{}^H_\AA\MMM_\AA^{\sym, \rm{fp}}}}
\def\dd{\mathrm{d}}
\def\ii{\mathrm{i}}
\def\oR{\bar R}
\def\pstd{{}^{\phantom{J}}_{\raisebox{\depth}{\scalebox{1.1}[-.75]{\mbox{${\mathbb
            \Delta}$}}}}}
\def\ra{\triangleright}
\def\stD{\dd\pstd}
\def\stDalpv{\stD{}_{\mbox{$\!{}^{}_{\oR^\al\ra v}\!$}}}
\def\stDalu{\stD{}_{\mbox{$\!{}^{}_{\oR_\al\ra u}\!$}}}
\def\stDav{\stD{}_{\mbox{$\!{}^{}_{{\,}^\al\!\;\! v}\!$}}}
\def\stDaz{\stD{}_{\mbox{$\!{}^{}_{{\:\!}^\al \!z}\!$}}}
\def\stDbaz{\stD{}_{\mbox{$\!{}^{}_{{\:\!}^{\beta\al}\! z}\!$}}}
\def\stDu{\stD{}_{\mbox{$\!{}^{}_u$}}}
\def\stDv{\stD{}_{\mbox{$\!{}^{}_v$}}}
\def\stDz{\stD{}_{\mbox{$\!{}^{}_z$}}}
\def\std{{\raisebox{\depth}{\scalebox{1.45}[-1]{\mbox{$\mathbb \Delta$}}}}_{\!}}
\def\stdd{\phantom{\!\!\!\dd}\pstd}
\def\sym{{\rm{sym}}}
\title{\bf Bianchi identities in noncommutative geometry}
\author{Paolo Aschieri}
\date{Source: arXiv:2608.24536v1 (CC BY 4.0)}
\begin{document}
\maketitle

\noindent\emph{Source and licence.} \bf Bianchi identities in noncommutative geometry. Version arXiv:2608.24536v1 (CC BY 4.0), distributed by arXiv under the Creative Commons Attribution 4.0 International licence, http://creativecommons.org/licenses/by/4.0/, as recorded in the arXiv licence metadata for this exact version and shown on the abstract page https://arxiv.org/abs/2608.24536v1. Full licence notice: \texttt{licenses/arxiv-2608.24536-CC-BY-4.0.txt}.\par\smallskip

\noindent\textit{Editorial scope.} Counted unit: the article's theorem BIANCHI2FORR (source0.tex lines 1381--1390). The article's complete original proof of it is delivered with it (source0.tex lines 1391--1425, one \texttt{proof} environment of 35 lines). No mathematics is written, repaired or re-derived: the statement and the proof are the article's own lines, in the article's own notation, and no retained line is substituted or re-flowed.  Retained uncounted as the source-local context the proof needs: lines 532--538; lines 1099--1102; lines 1117--1120.  Nothing was omitted from any retained span, so the package carries no omission record.  No retained line contains a citation, so the wrapper carries no bibliography and no bibliography entry.  No retained line contains a figure environment, an image-inclusion command or drawing code, so no image is delivered and the package declares no asset dependency.  Editorial formatting only, in addition: an AMS \texttt{amsart} preamble that restates the article's own declarations and macros used by the retained lines, namely the environments \texttt{theorem} on separate counters, together with the standard packages those lines need.  The retained environments print in this document as Theorem 1 in order of appearance, whereas the article prints the same items with the numbering of its own full document.  The complete native source of the article as distributed in the arXiv e-print for this exact version is bundled unchanged as \texttt{sources/208486/source0.tex}.
\begin{equation}\label{hexagon}
    (\Delta\otimes\mathrm{id})\mathcal{R}
    =\mathcal{R}_{13}\mathcal{R}_{23}
   ~,~~
    (\mathrm{id}\otimes\Delta)\mathcal{R}
    =\mathcal{R}_{13}\mathcal{R}_{12}
  \end{equation}

\begin{equation}\label{defCD}
\stD{}_{\mbox{$\!{}^{}_u$}}:=\ii_u\circ \stD+\stD
\circ \ii_u~,
\end{equation}

\begin{equation}\label{CRID}
\stDu\circ \ii_v -\ii_{\oR^\al\ra v} \circ \stDalu=\ii_{[u,v]}~~~~\mbox{
  i.e., }~~~~ \ii_u \circ \stDv -\stDalpv\circ \ii_{\oR_\al\ra u} =\ii_{[u,v]}~.
\end{equation}

\begin{theorem}\label{BIANCHI2FORR}
Let $\std$ be a left connection on an $A$-module $\Gamma\in\catfp$ and $R\stdd$ be its curvature tensor. The Bianchi identity for  $R\stdd$ is the vanishing of the linear maps $\Gamma\to \Gamma$ given by,
for all $u,v,z\in \Vect(A)$,
$$
\std{}_u\circ R\stdd(v,z,\mbox{-})-R\stdd(u,v,\mbox{-})\circ \std{}_z + R\stdd(u,[v,z],\mbox{-})+cp^{}_{\mathcal{R}}(u,v,z)=0
$$
or equivalently
$$
\std{}_u\circ R\stdd(v,z,\mbox{-})-R\stdd({}^\al v,{}^\be z,\mbox{-})\circ \std{}^{}_{{}_{\be\al} z} + R\stdd(u,[v,z],\mbox{-})+cp^{}_{\mathcal{R}}(u,v,z)=0
~.$$\end{theorem}

\begin{proof}
The Bianchi identity is equivalent to, for all $u,v,z\in \Vect(A)$, $\ii_u\circ \ii_v\circ \ii_z\circ \BDst({\stD}^{\!\!\!\!\!2}\:)=0$. Some preliminary observations are needed in order to compute this expression.   Recall definition \ref{defCD} in the form 
$\ii_z \stD=\stDz-\stD \ii_z$ (we omit the compositon $\circ$ of maps) and the Cartan relation in \eqref{CRID},
$\ii_v \stDz =\stDaz  {_{\,}}\ii_{{}_\al\! v} +\ii_{[v,z]}$.
Then move the inner derivatives to the right of the covariant derivative in the following expression to obtain 
$$
\ii_v\ii_z\stD=\stDaz  {_{\,}}\ii_{{}_\al v} +\ii_{[v,z]}-\stDv \ii_z+ \stD \ii_v\ii_z~.$$
Similarly,
\begin{equation}\label{iiid}\begin{split}
  \ii^{}_u \ii^{}_v\ii^{}_z\stD =~ & \stDbaz{}_{\,}\ii^{}_{{}^{}_\beta \!\!\; u\,} \ii^{}_{{}^{}_\al v}+ \ii^{}_{[u,{}^\alpha z]}\ii^{}_{{}_\alpha v}+ \ii^{}_u \ii^{}_{[v,z]} \\
  & -\stDav{} _{\,}\ii^{}_{{}^{}_\al \!\!\; u\,} \ii^{}_{z}- \ii^{}_{[u,v]}\ii^{}_{z}+
\stDu{}_{\,}\ii^{}_{v} \ii^{}_{z}-\stD  \ii^{}_u \ii^{}_v\ii^{}_z~.
\end{split}
\end{equation}
The last term in this expression drops out when acting on $\Omega^2(A)\otimes\Gamma$.
Next, similarly compute $\ii_u \ii_v \ii_z \stD\stD: \Omega(A)\otimes_A\Gamma\to  \Omega(A)\otimes_A\Gamma$ to be,
\begin{equation}\label{iiiddR}
\ii_u \ii_v \ii_z\stD\stD=- R\stdd(u,v,\mbox{-})\circ \ii_z +cp^{}_{\mathcal{R}}(u,v,z)~;
\end{equation}
to obtain this expression observe that six addends cancel each other out in pairs (use \eqref{hexagon}) and that the sum of three more addends vanishes due to the  ${\mathcal{R}}$-Jacoby identity.


For any $s\in \Gamma$, the left-hand side of the Bianchi identity   
$(\ii^{}_u \ii^{}_v\ii^{}_z\stD)\,{\stD}^{\!\!\!\!\!2\,}=(\ii^{}_u \ii^{}_v\ii^{}_z{\stD}^{\!\!\!\!\!2}\,)\,\stD$
reads, using \eqref{iiid},
\begin{equation}\label{RBL}
  (\ii^{}_u \ii^{}_v\ii^{}_z\stD)\,{\stD}^{\!\!\!\!\!2}\;s=-\std{}_u{\,}R\stdd(v,z,s)-R\stdd(u,[v,z],s)+cp^{}_{\mathcal{R}}(u,v,z)~.
  \end{equation}
The right-hand side reads 
\begin{equation}\label{RBR}
  (\ii^{}_u \ii^{}_v\ii^{}_z{\stD}^{\!\!\!\!\!2\,})\,\stD s=-R\stdd(u,v,{\std}{}_z s)+cp^{}_{\mathcal{R}}(u,v,z)~.
  \end{equation}
This proves the first identity of the theorem. The second identity  follows immediately from the equality
$u\otimes_Av\otimes_A z +cp^{}_{\mathcal{R}}(u,v,z)={}^\al v\otimes_A{}^\be z\otimes_A {}_{\be\al}u +cp^{}_{\mathcal{R}}(u,v,z)$ which holds since the triangular $R$-matrix induces a representation of the symmetric group $S_3$.
\end{proof}

\end{document}
